\section*{Bab \ref{ch:b2:continuous-reactors} --- Reaktor Kontinu dan Proses Industri}

\begin{solution}{exo:b2:continuous-reactors:1}
$\tau = V/Q_v = 2000/0.50 = \qty{4.0e3}{s}$, yaitu \qty{1.1}{h}.
\end{solution}

\begin{solution}{exo:b2:continuous-reactors:2}
$k\tau = 1.2 \times 10^{-3} \times 1800 = 2.16$. Tangki: $X = 2.16/3.16 = 0.68$;
pipa: $X = 1 - \mathrm e^{-2.16} = 0.88$.
\end{solution}

\begin{solution}{exo:b2:continuous-reactors:3}
Pipa: $\tau = \ln 100/0.020 = \qty{230}{s}$. Tangki: $X/(1 - X) = 99 = k\tau$,
$\tau = \qty{4950}{s}$, lebih dari dua puluh kali lebih lama: pada \omterm{def:b2:continuous-reactors:conversion}{konversi}
tinggi, tangki berpengaduk membayar mahal karena bekerja pada konsentrasi
keluaran.
\end{solution}

\begin{solution}{exo:b2:continuous-reactors:4}
Per liter: $120 \times 0.50 = \qty{60}{kJ}$ dilepaskan, memanaskan
\qty{4.2}{kJ/K}: $\Delta T_{\text{ad}} = \qty{14}{K}$. Kegagalan pendinginan
akan menghangatkan campuran secara nyata tetapi tidak berbahaya, kecuali
jika pelarutnya dekat dengan titik didihnya atau ada reaksi kedua yang
mulai berjalan dalam rentang itu.
\end{solution}

\begin{solution}{exo:b2:continuous-reactors:5}
Tangki: $Da = X/(1 - X)^2 = 0.8/0.04 = 20$, $\tau = 20/0.010 = \qty{2000}{s}$.
Pipa: $Da = X/(1 - X) = 4$, $\tau = \qty{400}{s}$. Rasionya 5, dibandingkan
dengan $4/\ln 5 =
2.5$ untuk orde 1: makin tinggi ordenya, makin besar
kerugian tangki berpengaduk karena konsentrasinya yang rendah.
\end{solution}

\begin{solution}{exo:b2:continuous-reactors:6}
$k\tau = 5.0$. Tiga tangki: $X = 1 - (1 + 5/3)^{-3} = 0.947$. Satu tangki:
$5/6 =
0.833$. Pipa: $1 - \mathrm e^{-5} = 0.993$.
\end{solution}

\begin{solution}{exo:b2:continuous-reactors:7}
$X = 0.70$ berarti $k\tau = 0.70/0.30 = 2.33$. Menggandakan laju alir
membagi dua $\tau$: $k\tau = 1.17$, $X = 0.54$. Untuk kembali mencapai
70\,\%, \omterm{def:b2:continuous-reactors:residence-time}{waktu tinggal}, dan dengan demikian volume, harus digandakan.
\end{solution}

\begin{solution}{exo:b2:continuous-reactors:8}
Dilepaskan: $80 \times 2.0 \times 1.0 \times 0.90 = \qty{144}{kW}$. Dibawa
oleh aliran keluaran, yang dipanaskan dari \qty{330}{K} ke \qty{350}{K}:
$4.0 \times 1.0 \times 20 =
\qty{80}{kW}$. Mantel harus membuang
selisihnya, \qty{64}{kW}.
\end{solution}

\begin{solution}{exo:b2:continuous-reactors:9}
Dimensi linear $\times 10$: volume $\times 1000$, luas dinding
$\times 100$, luas per satuan volume $\div 10$. Kalor yang dilepaskan
bertambah sebanding dengan volume, kalor yang dibuang melalui dinding
sebanding dengan luasnya: reaktor besar, per liter, membuang kalor sepuluh
kali lebih sedikit pada selisih suhu yang sama.
\end{solution}

\begin{solution}{exo:b2:continuous-reactors:10}
Pipa: $\tau = \ln(0.05/0.10)/(0.05 - 0.10) = \qty{13.9}{min}$, $c_B/c_{A,0} =
\frac{0.10}{-0.05}(\mathrm e^{-1.386} - \mathrm e^{-0.693}) = 0.50$. Tangki: $\tau =
1/\sqrt{0.005} = \qty{14.1}{min}$, $c_B/c_{A,0} = 1.414/(2.414 \times 1.707) =
0.34$. Di dalam tangki, A yang baru masuk dan B yang sudah jadi bercampur
dalam volume yang sama: sebagian B selalu tinggal cukup lama untuk diubah
menjadi C sementara sebagian A baru saja tiba; di dalam pipa, setiap
irisan mempunyai riwayat yang sama.
\end{solution}

\begin{solution}{exo:b2:continuous-reactors:11}
Menaikkan suhu pendingin (dan umpan) menggeser garis pembuangan ke kanan.
Titik potong dingin naik perlahan sampai garis itu menyinggung lengkungan
bawah kurva S; setelah itu keadaan dingin lenyap dan tangki melompat ke
cabang panas (penyalaan). Ketika suhu diturunkan lagi, tangki tetap
berada di cabang panas sampai garis itu menyinggung lengkungan atas, dan
baru kemudian jatuh kembali (pemadaman), pada suhu yang lebih rendah
daripada suhu penyalaan: sebuah lingkar histeresis.
\end{solution}

\begin{solution}{exo:b2:continuous-reactors:12}
Batch: tersisa 75\,\% dari \qty{2.0}{mol/L}, $1.5 \times 100 = \qty{150}{kJ/L}$, $\Delta T = 150/4.0 = \qty{37.5}{K}$ di atas suhu kerja,
cukup untuk mencapai titik didih atau suatu dekomposisi. Semi-batch:
paling banyak \qty{0.10}{mol/L} belum bereaksi, $\Delta T = 10/4.0 = \qty{2.5}{K}$. Penambahan perlahan membatasi energi yang tersimpan di
dalam reaktor.
\end{solution}

\begin{solution}{pb:b2:continuous-reactors:1}
\textbf{1.} $r = k[\text{ester}][\ce{OH-}]$; umpan yang sama dan
stoikiometri 1 : 1 menjaga kedua konsentrasi tetap sama, $r = kc^2$.
\textbf{2.} $1/c - 1/c_0 = kt$ dengan $c = 0.10c_0$: $kc_0t = 9$, $t = 9/(0.11
\times 0.10) = \qty{818}{s}$, sekitar \qty{14}{min}.
\textbf{3.} $t_{1/2} = 1/(kc_0) = \qty{91}{s}$; untuk orde 2 laju turun
sebanding dengan kuadrat konsentrasi, sehingga waktu untuk menjadikannya
separuh bergantung pada titik awalnya.
\textbf{4.} $Da = kc_0\tau = 0.011\,\tau$ ($\tau$ dalam detik).
\textbf{5.} $Q_vc_0 - Q_vc - kc^2V = 0$, yaitu $c_0 - c = k\tau c^2$.
\textbf{6.} Dengan $c = c_0(1 - X)$: $c_0X = k\tau c_0^2(1 - X)^2$.
\textbf{7.} $Da = 0.90/0.10^2 = 90$.
\textbf{8.} $\tau = 90/0.011 = \qty{8.2e3}{s}$, \qty{2.3}{h}.
\textbf{9.} $V = 0.50 \times 8182 = \qty{4.1e3}{L}$, \qty{4.1}{m^3}.
\textbf{10.} Konsentrasi keluaran, $0.010\,\unit{mol/L}$: lajunya di
mana-mana seratus kali lebih rendah daripada di saluran masuk.
\textbf{11.} $Da = X/(1 - X) = 9$, $\tau = \qty{818}{s}$, $V = \qty{409}{L}$.
\textbf{12.} Volume pipa sepuluh kali lebih kecil.
\textbf{13.} $c_{j-1} - c_j = k\tau_1c_j^2$, dengan $\tau_1$ \omterm{def:b2:continuous-reactors:residence-time}{waktu tinggal}
satu tangki.
\textbf{14.} $\tau_1 = 850/(3 \times 0.50) = \qty{567}{s}$, $k\tau_1 =
\qty{62.4}{L/mol}$. Penyelesaian $62.4c_j^2 + c_j - c_{j-1} = 0$:
$c_1 = 0.0328$, $c_2 = 0.0163$, $c_3 = \qty{0.0100}{mol/L}$: \omterm{def:b2:continuous-reactors:conversion}{konversi}
90\,\%.
\textbf{15.} Pipa, atau kaskade (sepuluh dan lima kali lebih kecil daripada
satu tangki). Satu tangki mungkin tetap unggul karena suhunya seragam,
pembersihan dan pengendaliannya mudah, dan karena di sini volume itu murah.
\textbf{16.} $\Delta T_{\text{ad}} = 55\,000 \times 100/(997 \times 4180) =
\qty{1.3}{K}$ (\qty{1.2}{K} pada 90\,\%).
\textbf{17.} $55 \times 0.50 \times 0.10 \times 0.90 = \qty{2.5}{kW}$.
\textbf{18.} Tidak: aliran keluaran membawa pergi kalornya dengan kenaikan
sekitar satu kelvin; \omterm{def:b2:continuous-reactors:runaway}{pelarian termal} tidak mungkin terjadi.
\textbf{19.} Dua puluh kali lebih pekat: kenaikan adiabatiknya menjadi
\qty{26}{K} dan kalornya \qty{50}{kW}, sehingga layak diberi koil
pendingin; karena laju $kc_0$ dua puluh kali lebih tinggi, tangki untuk
90\,\% akan dua puluh kali lebih kecil (\qty{0.20}{m^3}).
\textbf{20.} $0.50 \times 0.10 \times 0.90 = \qty{0.045}{mol/s}$, $\times 86\,400
\times 82.03 = \qty{319}{kg}$ natrium etanoat per hari.
\textbf{21.} Tetapan laju bertambah seiring suhu (Arrhenius), sehingga
volumenya menyusut; harganya adalah pemanasan umpan dan, untuk reaksi
lain, \omterm{def:b2:continuous-reactors:selectivity}{selektivitas} yang lebih rendah serta lebih banyak uap.
\textbf{22.} Satu tangki berpengaduk harus bervolume
$\boldsymbol{V \approx
\qty{4.1}{m^3}}$.
\end{solution}
